% GENERATED FILE. Edit canonical lessons, then run npm run generate:books.
% canonical-source-hash: fnv1a64:5d8daf0fb192a93e
% canonical-lessons: SA-C428-grhapatih, SA-C428-abhyagatah, SA-C428-nivasayati, SA-C428-pratyudgacchati, SA-C428-anuvrajati

\chapter{Hosts and Guests}
\label{ch:sa-hosts-and-guests}
\hlchaptermodality{428}

\begin{chapteropening}
\textbf{By the end of this chapter:} \emph{I can say how a host receives and sees off travellers.}

Complete the last lesson of chapter 428: I can say how a host receives and sees off travellers.
\end{chapteropening}

\section[\texorpdfstring{\sk{गृहपतिः} (gṛhapatiḥ)}{गृहपतिः (gṛhapatiḥ)}]{\sk{गृहपतिः} (gṛhapatiḥ) — the master of the house}

\label{lesson:SA-C428-grhapatih}



\begin{quote}
Before the new one: say the Sanskrit for spreads out, then the Sanskrit for wakes up.
\end{quote}

\subsection*{You'll want to know: \sk{गृहपतिः}}
\textbf{\sk{गृहपतिः}} — \emph{gṛhapatiḥ} — \textquotedblleft{}the master of the house\textquotedblright{}.

\begin{quote}
\sk{गृहपतिः} \sk{अतिथिम्} \sk{आह्वयति।}
\end{quote}

\begin{quote}
\emph{gṛhapatiḥ atithim āhvayati.}
\end{quote}

\begin{quote}
The master of the house calls in the guest.
\end{quote}

A masculine noun.

\begin{cousinweb}
\textbf{\sk{गृहपतिः}} (\emph{gṛhapatiḥ}) is built from \emph{gṛha}, \textquotedblleft{}a house\textquotedblright{}, and \emph{pati}, \textquotedblleft{}a lord, a master\textquotedblright{}; MW gives the compound as the master of a house, a householder (MW).
\end{cousinweb}

\begin{grammarlens}[title={What you've built}]
One more word for getting around.
\end{grammarlens}

\subsection*{Your turn}
\begin{itemize}
\raggedright
  \item \emph{Say it:} \emph{gṛhapatiḥ}
  \item \emph{Say it:} \emph{gṛhapatiḥ}, once more
  \item \emph{Say it:} \emph{gṛhapatiḥ atithim āhvayati.} — The master of the house calls in the guest.
  \item \emph{Recall:} say the Sanskrit for spreads out, then the Sanskrit for wakes up, then say \emph{gṛhapatiḥ} again
\end{itemize}

\begin{tcolorbox}[breakable,colback=teal!4,colframe=teal!35!black,title={Before you move on}]
What does \sk{गृहपतिः} mean? (\textquotedblleft{}The master of the house\textquotedblright{}.) Say it once more.
\end{tcolorbox}
\section[\texorpdfstring{\sk{अभ्यागतः} (abhyāgataḥ)}{अभ्यागतः (abhyāgataḥ)}]{\sk{अभ्यागतः} (abhyāgataḥ) — a visitor; a guest who calls}

\label{lesson:SA-C428-abhyagatah}



\begin{quote}
Before the new one: say the Sanskrit for wakes up, then the Sanskrit for the master of the house.
\end{quote}

\subsection*{You'll want to know: \sk{अभ्यागतः}}
\textbf{\sk{अभ्यागतः}} — \emph{abhyāgataḥ} — \textquotedblleft{}a visitor; a guest who calls\textquotedblright{}.

\begin{quote}
\sk{अभ्यागतः} \sk{द्वारे} \sk{तिष्ठति।}
\end{quote}

\begin{quote}
\emph{abhyāgataḥ dvāre tiṣṭhati.}
\end{quote}

\begin{quote}
The visitor stands at the door.
\end{quote}

A masculine noun.

\begin{cousinweb}
\textbf{\sk{अभ्यागतः}} (\emph{abhyāgataḥ}) is built from \emph{abhi}, \textquotedblleft{}towards\textquotedblright{}, \emph{ā}, \textquotedblleft{}near\textquotedblright{}, and the root √gam, \textquotedblleft{}to go\textquotedblright{}: one who has come up to the house; MW gives \emph{abhyāgata} as come, arrived and, as a noun, an uninvited guest or any guest (MW).
\end{cousinweb}

\begin{grammarlens}[title={What you've built}]
One more word for getting around.
\end{grammarlens}

\subsection*{Your turn}
\begin{itemize}
\raggedright
  \item \emph{Say it:} \emph{abhyāgataḥ}
  \item \emph{Say it:} \emph{abhyāgataḥ}, once more
  \item \emph{Say it:} \emph{abhyāgataḥ dvāre tiṣṭhati.} — The visitor stands at the door.
  \item \emph{Recall:} say the Sanskrit for wakes up, then the Sanskrit for the master of the house, then say \emph{abhyāgataḥ} again
\end{itemize}

\begin{tcolorbox}[breakable,colback=teal!4,colframe=teal!35!black,title={Before you move on}]
What does \sk{अभ्यागतः} mean? (\textquotedblleft{}A visitor; a guest who calls\textquotedblright{}.) Say it once more.
\end{tcolorbox}
\section[\texorpdfstring{\sk{निवासयति} (nivāsayati)}{निवासयति (nivāsayati)}]{\sk{निवासयति} (nivāsayati) — puts up for the night; lodges (a guest)}

\label{lesson:SA-C428-nivasayati}



\begin{quote}
Before the new one: say the Sanskrit for the master of the house, then the Sanskrit for a visitor.
\end{quote}

\subsection*{You'll want to know: \sk{निवासयति}}
\textbf{\sk{निवासयति}} — \emph{nivāsayati} — \textquotedblleft{}puts up for the night; lodges (a guest)\textquotedblright{}.

\begin{quote}
\sk{गृहपतिः} \sk{पथिकं} \sk{गृहे} \sk{निवासयति।}
\end{quote}

\begin{quote}
\emph{gṛhapatiḥ pathikaṃ gṛhe nivāsayati.}
\end{quote}

\begin{quote}
The master of the house puts the wayfarer up in his home.
\end{quote}

Past (he …): \sk{सः} \sk{पथिकं} \sk{गृहे} \sk{निवासितवान्}, \emph{saḥ pathikaṃ gṛhe nivāsitavān}.

\begin{cousinweb}
\textbf{\sk{निवासयति}} (\emph{nivāsayati}) is built from \emph{ni}, \textquotedblleft{}down, in\textquotedblright{}, and the causative of the root √vas, \textquotedblleft{}to dwell\textquotedblright{}, which gives \emph{vasati}, \textquotedblleft{}dwells\textquotedblright{}; MW gives \emph{ni-vāsayati} as to make stay, to receive as a guest, and KRM gives the causative participle \emph{vāsita} (MW).
\end{cousinweb}

\begin{grammarlens}[title={What you've built}]
One more word for getting around.
\end{grammarlens}

\subsection*{Your turn}
\begin{itemize}
\raggedright
  \item \emph{Say it:} \emph{nivāsayati}
  \item \emph{Say it:} \emph{nivāsayati}, once more
  \item \emph{Say it:} \emph{gṛhapatiḥ pathikaṃ gṛhe nivāsayati.} — The master of the house puts the wayfarer up in his home.
  \item \emph{Recall:} say the Sanskrit for the master of the house, then the Sanskrit for a visitor, then say \emph{nivāsayati} again
\end{itemize}

\begin{tcolorbox}[breakable,colback=teal!4,colframe=teal!35!black,title={Before you move on}]
What does \sk{निवासयति} mean? (\textquotedblleft{}Puts up for the night; lodges (a guest)\textquotedblright{}.) Say it once more.
\end{tcolorbox}
\section[\texorpdfstring{\sk{प्रत्युद्गच्छति} (pratyudgacchati)}{प्रत्युद्गच्छति (pratyudgacchati)}]{\sk{प्रत्युद्गच्छति} (pratyudgacchati) — goes out to meet}

\label{lesson:SA-C428-pratyudgacchati}



\begin{quote}
Before the new one: say the Sanskrit for a visitor, then the Sanskrit for puts up for the night.
\end{quote}

\subsection*{You'll want to know: \sk{प्रत्युद्गच्छति}}
\textbf{\sk{प्रत्युद्गच्छति}} — \emph{pratyudgacchati} — \textquotedblleft{}goes out to meet\textquotedblright{}.

\begin{quote}
\sk{गृहपतिः} \sk{अतिथिं} \sk{प्रत्युद्गच्छति।}
\end{quote}

\begin{quote}
\emph{gṛhapatiḥ atithiṃ pratyudgacchati.}
\end{quote}

\begin{quote}
The master of the house goes out to meet the guest.
\end{quote}

Past (he …): \sk{सः} \sk{अतिथिं} \sk{प्रत्युद्गतवान्}, \emph{saḥ atithiṃ pratyudgatavān}.

\begin{cousinweb}
\textbf{\sk{प्रत्युद्गच्छति}} (\emph{pratyudgacchati}) is built from \emph{prati}, \textquotedblleft{}towards\textquotedblright{}, \emph{ud}, \textquotedblleft{}up, out\textquotedblright{}, and the root √gam, \textquotedblleft{}to go\textquotedblright{}; MW gives \emph{praty-ud-gam} as to go out towards, to advance to meet a friend, with the participle \emph{pratyudgata} (MW).
\end{cousinweb}

\begin{grammarlens}[title={What you've built}]
One more word for getting around.
\end{grammarlens}

\subsection*{Your turn}
\begin{itemize}
\raggedright
  \item \emph{Say it:} \emph{pratyudgacchati}
  \item \emph{Say it:} \emph{pratyudgacchati}, once more
  \item \emph{Say it:} \emph{gṛhapatiḥ atithiṃ pratyudgacchati.} — The master of the house goes out to meet the guest.
  \item \emph{Recall:} say the Sanskrit for a visitor, then the Sanskrit for puts up for the night, then say \emph{pratyudgacchati} again
\end{itemize}

\begin{tcolorbox}[breakable,colback=teal!4,colframe=teal!35!black,title={Before you move on}]
What does \sk{प्रत्युद्गच्छति} mean? (\textquotedblleft{}Goes out to meet\textquotedblright{}.) Say it once more.
\end{tcolorbox}
\section[\texorpdfstring{\sk{अनुव्रजति} (anuvrajati)}{अनुव्रजति (anuvrajati)}]{\sk{अनुव्रजति} (anuvrajati) — sees a guest off}

\label{lesson:SA-C428-anuvrajati}



\begin{quote}
Before the new one: say the Sanskrit for puts up for the night, then the Sanskrit for goes out to meet.
\end{quote}

\subsection*{You'll want to know: \sk{अनुव्रजति}}
\textbf{\sk{अनुव्रजति}} — \emph{anuvrajati} — \textquotedblleft{}sees a guest off\textquotedblright{}.

\begin{quote}
\sk{गृहपतिः} \sk{अतिथिम्} \sk{अनुव्रजति।}
\end{quote}

\begin{quote}
\emph{gṛhapatiḥ atithim anuvrajati.}
\end{quote}

\begin{quote}
The master of the house sees the guest off.
\end{quote}

Past (he …): \sk{सः} \sk{अतिथिम्} \sk{अनुव्रजितवान्}, \emph{saḥ atithim anuvrajitavān}.

\begin{cousinweb}
\textbf{\sk{अनुव्रजति}} (\emph{anuvrajati}) is built from \emph{anu}, \textquotedblleft{}after\textquotedblright{}, and the root √vraj, \textquotedblleft{}to walk, to go\textquotedblright{}; MW gives \emph{anu-vraj} as to go along and to follow a departing guest some way as a mark of respect, and \emph{vrajita}, \textquotedblleft{}gone\textquotedblright{} (MW).
\end{cousinweb}

\begin{grammarlens}[title={What you've built}]
One more word for getting around.
\end{grammarlens}

\subsection*{Your turn}
\begin{itemize}
\raggedright
  \item \emph{Say it:} \emph{anuvrajati}
  \item \emph{Say it:} \emph{anuvrajati}, once more
  \item \emph{Say it:} \emph{gṛhapatiḥ atithim anuvrajati.} — The master of the house sees the guest off.
  \item \emph{Recall:} say the Sanskrit for puts up for the night, then the Sanskrit for goes out to meet, then say \emph{anuvrajati} again
\end{itemize}

\begin{tcolorbox}[breakable,colback=teal!4,colframe=teal!35!black,title={Before you move on}]
What does \sk{अनुव्रजति} mean? (\textquotedblleft{}Sees a guest off\textquotedblright{}.) Say it once more.
\end{tcolorbox}
